\documentclass[10pt]{article}
\usepackage[margin=20mm]{geometry}
\usepackage{amsmath,amssymb,mathtools,fontspec,unicode-math,xurl,hyperref}
\IfFontExistsTF{DejaVu Serif}{\setmainfont{DejaVu Serif}}{\setmainfont{Latin Modern Roman}}
\setmathfont{latinmodern-math.otf}
\setlength{\parindent}{0pt}\setlength{\parskip}{6pt}\sloppy
\title{Local densities and lattice indices in the proof of large prime factors of quartic polynomial values}
\author{AI-assisted technical review}
\date{}
\begin{document}
\maketitle
Abstract: The 86-page accepted manuscript of Dartyge and Maynard is examined through its algebraic reduction, sieve parameters, and Type I/II estimates. A joint local-density calculation gives a counterexample to the constant in Lemma 9.5. The Type II argument also uses an invalid uniform lattice-index identification and an unsupported truncation of a minors gcd. These are Major issues in the presented proof, without a counterexample to the principal theorem. Two further Major proof defects concern the exponent saving and the prime reciprocal sum in the main Type I term. A Minor group concerns corrected discriminant identities and a restricted root-count boundary. Four Major groups and one Minor group are established. Explicit computations and the limits of the proposed repairs are given. Scientific adjudication remains pending.\par
\section*{Identity, scope, and proof structure}
The reviewed work is Cécile Dartyge and James Maynard, On the largest prime factor of quartic polynomial values: the cyclic and dihedral cases, DOI 10.4171/JEMS/1586. The scientific source is the 86-page Oxford Research Archive accepted author manuscript [paper], with repository acceptance date 12 October 2023 and deposit date 24 October 2023. The PDF’s internal creation date, 12 September 2023, is recorded separately. The journal page gives 31 January 2025 as publication date and restricts the journal PDF to subscribers [journal]. The acquired 84-page arXiv v1 of 7 December 2022 is a distinct version [arxiv]; neither its TeX nor its extracted text is treated as equivalent to the AAM.\par
During the initial review, the complete AAM text, including definitions, hypotheses, all proof sections, and references, was read. All source pages cited as decisive evidence below were also viewed as rendered images. The review checks the internal proof chain; it does not independently re-prove all results imported from the literature, in particular the full arguments of reference [14]. Source and report page receipts record exactly which images were viewed. All page numbers below refer to one-based AAM PDF pages, which match the printed numbers.\par
Theorem 1.1 (p. 3) asserts a positive-proportion lower bound \(P^+(P(n))\ge x^{1+c_P}\) on \(x<n\le2x\) for a monic irreducible quartic with Galois group \(C_4\) or \(D_4\). Lemma 3.1 and Propositions 3.2–3.3 (pp. 7–8) reduce this to \(S_0\gg1\) and \(S_1=o(X)\). Theorem 4.1 (pp. 9–11) is the central localized incomplete-norm-form estimate. Section 5 expresses a sextic auxiliary form as a quartic norm factor \(q_1\) times a quadratic factor \(q_2\); Sections 6–8 use its controlled factorization. Sections 9–11 supply the Type I/II estimates. The findings below distinguish failures of auxiliary assertions from failures of the final theorem.\par
\section*{Joint local densities and the Type I constant}
Problem 1 (Major: the constant in Lemma 9.5 is incorrect as stated).  Equation (9.3) and Lemma 9.4 are on p. 50; the Euler product (9.4) and Lemma 9.5 are on p. 51. The multiplicativity proved there assumes coprime norms. The ensuing product over individual prime ideals does not respect this restriction.\par
For an integral ideal \(\mathfrak d\), the definition is \[\rho_v(\mathfrak d)=\mathrm N(\mathfrak d)^{-2}
 \#\{(a_1,a_2,a_3)\bmod\mathrm N(\mathfrak d):
 \mathfrak d\mid(a_1\nu_1+a_2\nu_2+a_3\nu_3)\}.\] Take \(K=\mathbb Q(\zeta_5)\), \(v=(1,\zeta_5,\zeta_5^2,\zeta_5^3)\), \(m=1\), and \(I=\mathcal O_K\). These satisfy the stated quartic and basis hypotheses. At \(p=11\), the polynomial \(\Phi_5(T)=T^4+T^3+T^2+T+1\) has roots \(3,4,5,9\), and \((11)=\mathfrak p_1\mathfrak p_2\mathfrak p_3\mathfrak p_4\). Each degree-one factor has \(\rho_v(\mathfrak p_i)=1\), whereas \[\rho_v((11))=11,\] because divisibility by \((11)\) imposes divisibility of all three coefficients by 11. The normalization uses \(\mathrm N((11))=11^4\). This is not a counterexample to Lemma 9.4 itself: the norms of the four prime ideals are not coprime.\par
For a subset of \(k\) distinct roots, the simultaneous evaluation map on polynomials of degree at most two has rank \(\min(k,3)\). Direct enumeration of all \(11^3\) triples independently gives \[\begin{array}{c|rrrrr}
 k&0&1&2&3&4\\ \hline
 \text{solutions for each subset}&1331&121&11&1&1.
\end{array}\] Thus the actual squarefree Möbius local factor at \(s=1\) is \[A_p=1-\frac4p+\frac6{p^2}-\frac3{p^3},\] whereas independent multiplication over the four prime ideals gives \((1-p^{-1})^4\). At \(p=11\), these are respectively \(10010/14641\) and \(10000/14641\), with ratio \(1001/1000\).\par
\subsection*{A global contradiction, not just a local mismatch}
Put \[F(s)=\sum_{\mathfrak d}
 \frac{\mu(\mathfrak d)\rho_v(\mathfrak d)}{\mathrm N(\mathfrak d)^s}
 =\zeta_K(s)^{-1}G(s).\] The legitimate Euler grouping is by rational primes. Write \(\widetilde{\mathfrak S}\) for the product in (9.4). For \(p\ne5\), the splitting types in this field give the following ratios between the correct local factor of \(G(1)\) and the paper’s factor: \[\begin{array}{c|c}
 p\bmod5&\text{ratio}\\ \hline
 1&1+(p-1)^{-3}\\
 4&1+(p-1)/(p^2-1)^2\\
 2,3&1 .
 \end{array}\] For the type \(2+2\), divisibility by either degree-two prime imposes two conditions and divisibility by both imposes only three. The actual factor is \(1-2p^{-2}+p^{-3}\), yielding the displayed ratio. For an inert prime and for \(p=5\) there is only one prime ideal, so no cross-prime-ideal factor is missing. These extra factors form an absolutely convergent product, all are at least one, and the factor at 11 is strictly larger than one. Consequently \[G(1)>\widetilde{\mathfrak S}.\] The Euler product after cancellation of its degree-one terms converges in a neighborhood of \(s=1\); this justifies the value and continuity used here.\par
Define the sum in Lemma 9.5 for the chosen \(m,I\) by \[S(R)=\sum_{\mathrm N(\mathfrak d)<R}
 \frac{\mu(\mathfrak d)\rho_v(\mathfrak d)}{\mathrm N(\mathfrak d)}
 \log\frac R{\mathrm N(\mathfrak d)}.\] For \(u>0\), termwise integration in the absolutely convergent half-plane gives \[F(1+u)=u^2\int_1^\infty S(R)R^{-u-1}\,dR .\] If Lemma 9.5 held, \(S(R)\to\widetilde{\mathfrak S}/\gamma_K\), so the Abelian identity would imply \(\lim_{u\downarrow0}F(1+u)/u=\widetilde{\mathfrak S}/\gamma_K\). The simple pole of \(\zeta_K\), however, gives \[\lim_{u\downarrow0}\frac{F(1+u)}u=\frac{G(1)}{\gamma_K}
 >\frac{\widetilde{\mathfrak S}}{\gamma_K},\] a contradiction. No unproved replacement Tauberian theorem is needed for this refutation of the stated constant.\par
\subsection*{Impact and repair boundary}
Lemma 9.5 supplies the normalization for the approximation weights, and the same individual-prime-ideal separation appears in (10.17)–(10.20), pp. 64–66. This is a central proof dependency, so the supported severity is Major. The corrected joint local constant is explicit for \(m=1,I=\mathcal O_K\). A full repair must also establish the uniform formula for arbitrary \(m,I\), the associated \(h(J)\), and the Type I main term with the joint local factors. A multivariable local telescoping argument may retain the intended cancellation, but that complete repair is not proved here. Theorem 4.1 and Theorem 1.1 are not thereby disproved.\par
\section*{The Type II lattice argument}
Problem 2 (Major: restriction of congruences and a gcd truncation).  The proof of Lemma 11.7 on p. 76 identifies \(f_{\mathbf b_1,\mathbf b_2,p_1,p_2}=p_1p_2\) when \(p_1\ne p_2\). The justification proves nontriviality on the ambient four-dimensional space, not on the rank-two lattice to which the conditions are restricted. The preceding Lemmas 11.5–11.6, pp. 73–75, also use an invalid \(D_{\mathbf b_1,\mathbf b_2}\le B\) assertion.\par
\subsection*{An exact restricted-index counterexample}
Use \(K=\mathbb Q(\theta)\), \(\theta^4=-2\), with \(v=(1,\theta,\theta^2,\theta^3)\), and \(f(x)=x_1^2-5x_2^2\), \(D_f=20\). For \(b=(b_1,b_2,b_3,b_4)\), the fourth row of multiplication by \(b\) is \[T(b)=(b_4,b_3,b_2,b_1).\] Let \(p_1=41,p_2=61,s_1=13,s_2=26\), so \(s_i^2=5\bmod p_i\). The following explicit primitive pair was checked with integer arithmetic: \[\mathbf b_1=(250100001,1,0,0),\qquad
 \mathbf b_2=(250100367,-2501,3416,3294).\] Its norm residues modulo \(41,61\) are respectively \((3,3)\) and \((37,1)\); the two multiplication matrices are invertible modulo both primes. Writing \(M_i\) for these matrices, the two-row matrix \(T=(T(\mathbf b_1);T(\mathbf b_2))\) has rank two modulo each prime. The ranks after adjoining the relevant row \((1,s_i,0,0)M_i\) are \[\begin{array}{c|cc}
 &\operatorname{rank}T&\operatorname{rank}(T;(1,s_i,0,0)M_i)\\ \hline
 p_1=41&2&2\\
 p_2=61&2&3.
 \end{array}\] Since the rank-two row matrix has a unit minor at each prime, its integer kernel reduces onto the modular kernel. The first condition therefore has index one on \(\Lambda_{\mathbf b_1,\mathbf b_2}\), and the second has index 61. Hence the actual combined index is \[f_{\mathbf b_1,\mathbf b_2,41,61}=61\ne2501.\]\par
This phenomenon persists with growing primes. Choose distinct \(p_i\equiv1\pmod{20}\), put \(Q=p_1p_2\), take \(N\equiv1\pmod Q\), and let \(\mathbf b_1=(N,1,0,0)\). Choose \(\mathbf b_2=(M,-CQ,c_3,c_4)\), with fixed positive integer \(C\), so \[\mathbf b_2\equiv(-2,0,s_1,1+s_1)\pmod{p_1},\qquad
 \mathbf b_2\equiv(1,0,0,0)\pmod{p_2},\] \(M=N+O(Q)\), and \(c_3,c_4\in[Q,2Q)\). CRT permits these choices and, if needed, an additional condition \(\gcd(M,C)=1\); then the pair is primitive. The first new row equals \(T(\mathbf b_2)\bmod p_1\). Modulo \(p_2\), the original equations are \(a_3+a_4=a_4=0\), and the new one is \(a_1+s_2a_2=0\). Again the index is \(p_2\). The norm of \[-2+s\theta^2+(1+s)\theta^3
 =(s\theta^2+\theta^3)(1+\theta)\] is \(12(s^4+2)=324\bmod p_1\); thus the requisite norm coprimality holds.\par
The scale condition \(B\gg Q\) with a power margin follows from (4.11), since \(B=X^{S/4+o(1)}\), \(1<S<2\), and \(Q\le X^{2\tau'}\). The vectors approach the positive scalar direction, have norms \(N^4(1+O(Q/N))\), and can lie in a common permitted small box and a unit fundamental domain containing that direction. Choosing the sign \(-CQ\) also makes the real kernel compatible with positive first three product coordinates: \[a_4=-a_3/N,\qquad c_4a_1+c_3a_2=(CQ+M/N)a_3.\] Take \(C\) sufficiently large and \(a_2\asymp a_1/C\). The kernel equations give \(a_3\asymp a_1/C\) and \(a_4=-a_3/N\), so the \(\alpha\)-direction can also lie near the positive scalar ray. These observations exclude a mere fixed-small-prime or empty-real-region explanation for the local index error.\par
The calculation does not prove that the specific sieve weights \(g_{\mathbf b_1}g_{\mathbf b_2}\) are nonzero on this family, or that their total weighted contribution violates the final estimate. It refutes the uniform index assertion used to replace the summand. An exceptional-pair estimate or a different argument is required.\par
\subsection*{The minors gcd need not be at most \(B\)}
For the same power basis, take \[\mathbf b_1=(N,1,0,0),\qquad \mathbf b_2=(1,N-1,0,0),\qquad
 B=(N^4+2)^{1/4}.\] For \(N\ge7\), both norms \(N^4+2\) and \(1+2(N-1)^4\) lie in \([B^4,2B^4]\). Both vectors are primitive, but \[D_{\mathbf b_1,\mathbf b_2}=N(N-1)-1\asymp B^2.\] The common kernel is \(\mathbb Ze_1+\mathbb Ze_2\), of determinant one. At \(N=100\), \(D=9899\), and the vectors are nonproportional modulo 7. This directly invalidates the inference from individual primitivity to \(D\le B\) on p. 73. The ensuing sums truncate at \(d\le B\) on pp. 74–75 and again in the later reduction. For the example \(\|\mathbf b_1\wedge\mathbf b_2\|\asymp B^2\); the related determinant bound based on \(D\le B\) fails as well.\par
This is a counterexample to the intermediate geometric implication, not to the averaged bounds of Lemma 11.5 or 11.6. Membership of this second family in every particular fixed fundamental domain and in the full weight support is not asserted. A proof excluding or estimating \(D>B\) is still needed; simply extending the summation to \(B^2\) does not preserve the displayed error estimate. Together these issues leave the Type II proof incomplete as written.\par
\section*{Further proof corrections}
Problem 3 (Major: the exponent saving uses the wrong bound for \(q_{22}\)).  On p. 34, the bound for \(\theta_{22}\) is \((1+\alpha_0)/2-\theta_{21}\), but the second final condition substitutes \((1+\alpha_0)/4-\theta_{21}\).\par
The corrected condition is \[\theta_{21}>6\theta_0+10\alpha_0+
 512(2\eta_0+5\alpha_0/4).\] The p. 28 example has \(\alpha_0=10^{-5}\), \(\theta_{21}=0.001\); with \(\eta_0>\alpha_0\) its right side exceeds \(0.01674\). Thus that example cannot establish the required saving. A restricted, explicitly verified replacement is \[\alpha_0=10^{-7},\quad \theta_0=\tau_{ij}=10^{-10},
 \quad\eta_0=1.01\cdot10^{-7},\] with the original values \[(\theta_{11},\ldots,\theta_{16},\theta_{21})
 =(.1398,.1401,.1402,.21,.19,.1799,.001).\] Exact rational calculations verify (6.2)–(6.9), Lemma 7.1, the corrected three concluding conditions of Section 7, and the displayed condition of Lemma 8.1. The second corrected condition has margin \(0.0008315754\); (6.5) has margin \(4.94\cdot10^{-8}\). These positive margins allow the small epsilon losses in that calculation. This repair covers those explicit conditions, not every implicit choice later in the proof, and does not repair Problems 1–2.\par
Severity and actual dependency. The erroneous substitution is used on p. 34 to conclude \(S_1=o(X)\), the estimate required by Proposition 3.3. It therefore interrupts a core proof step, rather than merely misreporting an optional example. The explicit parameter repair is useful counterevidence to general infeasibility, but does not reduce the severity of the original derivation.\par
Problem 4 (Minor: the two representations of the discriminant factor).  Lemma 5.14 on p. 25 has an incorrect \(a_2^2\) coefficient and an incorrect representation in terms of \(q_2,q_3\).\par
An independent companion-matrix computation gives \(\Delta_{14}=-q_3h\), where \[\begin{aligned}
h={}&-4a_1^2+4c_3a_1a_2+(4c_2-3c_3^2)a_1a_3-c_3^2a_2^2\\
 &+(c_3^3-4c_1)a_2a_3+
 (-c_2c_3^2+4c_1c_3-4c_0)a_3^2\\
={}&-(c_3^2-4t_2)q_3-4q_2,\qquad
 t_2=(r_1+r_2)(r_3+r_4).\end{aligned}\] Thus \(-c_3a_2^2\) must be \(-c_3^2a_2^2\), and the last line of the printed identity must have the indicated sign and factor four. The exact symbolic differences are zero. Crucially, substitution of a root of \(q_1\) into this corrected \(h\) still gives the factorization (5.29). For odd primes satisfying the full coprimality and squarefreeness hypotheses of Lemma 5.16, the corrected formula retains the zero tests used there; the degenerate \(q_2\) branch is addressed separately below. This limits the impact of the displayed algebra error.\par
A related boundary deserves separate care. For the allowed \(D_4\) polynomial \(P(T)=T^4+2\), choose \(r_2=-r_1,r_3=ir_1,r_4=-r_3\). Then \(t_1=t_2=0\), contrary to Remark (ii) on p. 26. No subsequent division by \(t_1\) was found. More consequentially, \(c_3^2-4t_2=0\), so Lemma 5.16(ii) gives one \(B_{14}\) root at the \(q_2\)-primes, whereas the CRT calculation on p. 36 uniformly substitutes two. Increasing \(q_0\) cannot exclude divisors of zero. The accurate local factor uses \(\kappa_p\in\{1,2\}\) and \(1-\kappa_p/p\), which is at least \(1-2/p\). This suggests a lower-bound repair for the application, but the full weighted error analysis is not established here. No tuple satisfying every asymptotic condition of \(\mathcal J\) is claimed as a counterexample to the entire formula of Lemma 8.1.\par
Problem 5 (Major: normalization of the prime reciprocal sum in Type I).  The actual image of (10.22), p. 66, has \(\log(P_2/P_1)\). The corresponding main terms on pp. 59 and 61 carry the same expression.\par
With \(P_1=X^\tau,P_2=X^{\tau'}\), the prime reciprocal sum has main term \[\frac2{\varphi(D_f)}
 \log\frac{\log P_2}{\log P_1}
 =\frac2{\varphi(D_f)}\log\frac{\tau'}{\tau},\] rather than \(2\log(P_2/P_1)/\varphi(D_f)\). This follows by partial summation from the prime number theorem in the fixed progression. Theorem 4.1 itself has the correct \(\log(\tau'/\tau)\), so this is an intermediate normalization error, not a counterexample to its final constant.\par
Severity and actual dependency. The printed quantity is \((\tau'-\tau)\log X\), whereas the correct main term is bounded for fixed \(\tau,\tau'\). This is an error of asymptotic order in (10.22), propagated through Lemma 10.2 (p. 64), Proposition 9.13 (p. 59), and (9.28)–(9.30). The proof on pp. 59–60 explicitly derives Theorem 4.1 from Propositions 9.13 and 9.14. The already correct final logarithm is important counterevidence to a false final constant; it does not validate the printed Type I derivation. The core proof defect is Major even though partial summation identifies the necessary replacement.\par
There is also a finite-factor bookkeeping issue in (10.19)–(10.20): the ideals \(I\) in (9.20) exclude primes above \(m\). Formally writing \(r_e=\rho_v(\mathfrak P^e)/\mathrm N(\mathfrak P^e)\), \[h(\mathfrak P^e)=\frac{r_e-r_{e+1}}{1-r_1},\qquad
 1+\sum_{e\ge1}h(\mathfrak P^e)=\frac1{1-r_1}.\] Deleting factors above \(m\) cancels the earlier \(g((m'))\), since \(m'=mm_0\) and \(m_0\mid m^\infty\). This explains a possible cancellation of that finite factor, but it is only prime-local bookkeeping: Problem 1 prevents interpreting the formal calculation as a completed global repair.\par
\section*{Counterevidence, remaining scope, and assessment}
The general proof strategy is coherent at several important points. The quartic cofactor relation eliminates the \(a_0\)-dependence of the denominator; the factor \(q_1\) is represented as an incomplete quartic norm; the coprimality and squarefree hypotheses are used in the resultant argument. The corrected discriminant retains (5.29), and the smaller parameter package meets the explicit Section 7 constraints. The final scale requirements on p. 84, \[A<B^{3-\varepsilon},\qquad BP_2^2<A^{1-\varepsilon},\qquad
 P_2^{21}<B^{1-\varepsilon},\] are consistent with (4.10)–(4.11). These checks are counterevidence to an allegation of a general contradiction in the setup.\par
Item | Established result | Repair status\par
Problem 1 | Major; Lemma 9.5’s constant is refuted | Joint factors needed throughout Type I\par
Problem 2 | Major; uniform index and gcd implication fail | Exceptional-pair and large-gcd bounds remain\par
Problem 3 | Major; core exponent saving is unsupported | Replacement explicit inequalities verified\par
Problem 4 | Minor; discriminant displays are wrong | Correct identity and zero-test compatibility verified\par
Problem 5 | Major; Type I main-term order is wrong | Correct reciprocal-prime main term explicit\par
The revised assessment is Major, with four Major groups and one Minor group. This classification concerns substantive gaps in indispensable proof dependencies; it does not require that the paper be false or irreparable. No counterexample to the principal theorem is established, and no complete repair of the Type I/II argument is claimed. Scientific decision and adjudication remain pending, with no human reviews supplied.\par
Focused independent acceptance checked all five existing groups against the exact AAM, including the actual Type I/II dependencies and the limits of each counterexample. Problems 3 and 5 are raised from Minor to Major because their erroneous steps are used in core derivations; repairability is assessed separately. No new whole-paper initial review or independent proof of every imported result is claimed.\par
The bounded source search found no separate supplement or author code; absence of code is not a defect. The inspected current arXiv record lists only v1. No correction resolving these findings was identified on the checked primary records. Applicability to the subscription-only journal-final text remains unverified.\par
9 C. Dartyge and J. Maynard, On the largest prime factor of quartic polynomial values: the cyclic and dihedral cases, accepted manuscript, 86 pp., Oxford Research Archive, deposited 24 October 2023. https://ora.ox.ac.uk/objects/uuid:6d54ad71-14f1-42ba-a89a-3892ad6870f2. EMS Press, article record for DOI 10.4171/JEMS/1586, consulted 26 October 2026. https://ems.press/journals/jems/articles/14298551. C. Dartyge and J. Maynard, arXiv:2212.03381v1, 7 December 2022. https://arxiv.org/abs/2212.03381.\par
\section*{Evidence and attribution}
The exact reviewed PDF SHA-256 is\par
`d2b41d349d72b99db6d40f933d000717abd3d4a7068de0983080e222e22317e0`.\par
During the initial review, the root executor read the complete AAM, developed and verified the reported mathematical arguments, ran independent exact integer, finite-field, rational, and symbolic checks, compiled this report with the existing approved compiler, and inspected its final pages. The executable check source, standard output, standard error, exit code, and before/after hashes are retained with the evidence.\par
Two native helpers in the initial-review conversation supplied bounded cross-checks: `/root/check\_0211` covered algebra, parameters, local densities, and selected proof dependencies; `/root/check\_0213\_dispersion` covered the Type II lattice and dispersion chain. Neither performed compilation, PDF-image inspection, or independent acceptance. Their model backends were not independently verified. The controller’s UI observation is attributed to the controller; the host self-description is separately recorded. No equivalence between these observations is inferred. Native acceptance assistance supplied a bounded read-only challenge; it did not perform these executions, establish backend identity, or supply human review.\par
\end{document}
